\documentclass[11pt]{article}

\usepackage[margin=1in]{geometry}
\usepackage{amsmath, amssymb, amsthm}
\usepackage{booktabs}
\usepackage{array}
\usepackage{enumitem}
\usepackage{hyperref}

\newtheorem{proposition}{Proposition}
\newtheorem{assumption}{Assumption}
\theoremstyle{remark}
\newtheorem{remark}{Remark}

\newcommand{\R}{\mathbf{R}}
\newcommand{\xv}{\mathbf{x}}
\newcommand{\bt}{\boldsymbol{\beta}}
\newcommand{\dhat}{\hat{\mathbf{d}}}
\newcommand{\av}{\mathbf{a}}
\newcommand{\grad}{\nabla}
\newcommand{\dive}{\grad\!\cdot\!}
\newcommand{\Rd}{\mathbb{R}^d}

\title{\textbf{Schr\"odinger Bridges with a Reference Drift,\\
and an Analytic Shock-Rotation Drift for a Diamond Airfoil}}
\author{}
\date{}

\begin{document}
\maketitle

\begin{abstract}
We formulate the Schr\"odinger-bridge problem as a constrained minimization of
kinetic energy, derive its optimality system and Hopf--Cole linearization, and
express the solution through log-potentials and the iterative proportional
fitting procedure (IPFP). We then extend the formulation to a reference diffusion
carrying a macroscopic drift $\bt$, obtaining transport--diffusion equations for
the Hopf--Cole potentials. Finally we construct an explicit drift $\bt$ for the
compressible flow past a diamond airfoil, built from the oblique-shock geometry,
that transports the shock structure between two operating conditions.
\end{abstract}

\tableofcontents

% =====================================================================
\section{Notation}

Vectors lie in $\Rd$ (with $d=2$ in the application); $\|\cdot\|$ is the
Euclidean norm. The symbol $\gamma$ denotes the diffusion (entropic
regularization) coefficient of the bridge; in Section~\ref{sec:diamond} the
ratio of specific heats of the gas is written $\gamma_g$ to keep the two
distinct.

\begin{center}
\renewcommand{\arraystretch}{1.2}
\begin{tabular}{@{}>{$}l<{$}p{0.76\linewidth}@{}}
\toprule
\text{Symbol} & Meaning \\
\midrule
\mu,\nu & initial and terminal marginals, $\mu,\nu\in\mathcal{P}(\Rd)$ \\
t & time, rescaled to $[0,1]$ \\
\rho_t(\xv) & density at time $t$; $\rho_0=\mu$, $\rho_1=\nu$ \\
b_t(\xv) & control (velocity) field \\
\gamma & diffusion coefficient of the bridge \\
\mathcal{J} & kinetic-energy functional \\
\phi_t(\xv) & Lagrange multiplier / value function \\
\eta_t,\ \eta^\ast_t & forward / backward Hopf--Cole wave functions, $\rho_t=\eta_t\eta^\ast_t$ \\
f_t=\log\eta_t,\ g_t=\log\eta^\ast_t & log-potentials \\
\mathcal{P}_\tau & heat semigroup $e^{\tau\gamma\Delta}$ (Neumann BC) \\
\bt_t(\xv) & reference drift of the underlying diffusion \\
\mathcal{Q}_\tau,\ \mathcal{Q}_\tau^{\dagger} & reference transition operator and its adjoint \\
\alpha & wedge half-angle of the diamond \\
\gamma_g & ratio of specific heats of the gas \\
M & freestream Mach number; $M_0,M_1$ its endpoint values \\
\theta_{\mathrm{LE}},\theta_{\mathrm{TE}} & flow-deflection angles at leading / trailing edge \\
\beta_s & oblique-shock wave angle \\
\nu(\cdot) & Prandtl--Meyer function \\
\av_{\mathrm{LE}},\av_{\mathrm{TE}} & leading- / trailing-edge apex (pivot) points \\
\sigma\in\{+1,-1\} & surface index (upper / lower) \\
\kappa\in\{\mathrm{LE},\mathrm{TE}\} & shock-family index \\
\R & the rotation matrix $\left(\begin{smallmatrix}0&-1\\1&0\end{smallmatrix}\right)$ \\
M_0(\xv),M_1(\xv) & endpoint Mach fields on the mesh (the quantities interpolated) \\
T,\ S & forward ($\mu\!\to\!\nu$) and backward ($\nu\!\to\!\mu$) barycentric transport maps \\
\widehat M(t,\xv) & CDI-reconstructed Mach field at time $t$ \\
\mathrm{id}(\xv)=\xv & the identity (coordinate) map on the mesh \\
\bottomrule
\end{tabular}
\end{center}

% =====================================================================
\section{The Schr\"odinger problem}
\label{sec:sb}

Let $\mu,\nu\in\mathcal{P}(\Rd)$ be two probability measures and $\gamma>0$. The
dynamic Schr\"odinger problem seeks a density--velocity pair
$(\rho_t,b_t)_{t\in[0,1]}$ minimizing the expected kinetic energy
\begin{equation}
\label{eq:J}
\mathcal{J}(\rho,b)=\int_0^1\!\!\int_{\Rd}\tfrac12\,\|b_t(\xv)\|^2\,\rho_t(\xv)\,d\xv\,dt,
\end{equation}
subject to the Fokker--Planck constraint and the temporal boundary conditions
\begin{equation}
\label{eq:FP}
\partial_t\rho_t + \dive(\rho_t b_t) = \gamma\,\Delta\rho_t,
\qquad \rho_0=\mu,\quad \rho_1=\nu.
\end{equation}
Equivalently, $(\rho,b)$ is the law of the controlled diffusion
$dX_t=b_t(X_t)\,dt+\sqrt{2\gamma}\,dW_t$ that is closest, in relative entropy, to
the reference Brownian motion of the same intensity; by Girsanov's theorem the
relative entropy equals $\mathcal{J}/(2\gamma)$.

% =====================================================================
\section{Optimality system and linearization}
\label{sec:opt}

\subsection{Euler--Lagrange equations}
Introduce the multiplier $\phi_t(\xv)$ for the constraint \eqref{eq:FP} and form
\begin{equation}
\label{eq:Lag}
\mathcal{L}(\rho,b,\phi)=\int_0^1\!\!\int_{\Rd}\tfrac12\|b_t\|^2\rho_t\,d\xv\,dt
+\int_0^1\!\!\int_{\Rd}\phi_t\bigl(\partial_t\rho_t+\dive(\rho_t b_t)-\gamma\Delta\rho_t\bigr)\,d\xv\,dt.
\end{equation}

\paragraph{Variation in $b$.}
\begin{equation}
\delta_b\mathcal{L}=\int_0^1\!\!\int_{\Rd}\rho_t\,b_t\cdot\delta b_t\,d\xv\,dt
+\int_0^1\!\!\int_{\Rd}\phi_t\,\dive(\rho_t\,\delta b_t)\,d\xv\,dt .
\end{equation}
Integrating the second term by parts (no-flux boundary), $\int\phi\,\dive(\rho\,\delta b)=-\int\rho\,\grad\phi\cdot\delta b$, gives $\int_0^1\!\!\int\rho_t(b_t-\grad\phi_t)\cdot\delta b_t=0$. Since $\rho_t>0$ and $\delta b_t$ is arbitrary,
\begin{equation}
\label{eq:bopt}
b_t=\grad\phi_t .
\end{equation}

\paragraph{Variation in $\rho$ (with $\delta\rho_0=\delta\rho_1=0$).}
Integrating by parts in time, in the advection term, and twice in the diffusion term (self-adjointness of $\Delta$),
\begin{equation}
\delta_\rho\mathcal{L}=\int_0^1\!\!\int_{\Rd}\Bigl(\tfrac12\|b_t\|^2-\partial_t\phi_t-b_t\cdot\grad\phi_t-\gamma\Delta\phi_t\Bigr)\delta\rho_t\,d\xv\,dt=0,
\end{equation}
whence $\partial_t\phi_t+b_t\cdot\grad\phi_t+\gamma\Delta\phi_t-\tfrac12\|b_t\|^2=0$. Substituting \eqref{eq:bopt} yields the Hamilton--Jacobi--Bellman equation
\begin{equation}
\label{eq:HJB}
\partial_t\phi_t+\tfrac12\|\grad\phi_t\|^2+\gamma\,\Delta\phi_t=0 .
\end{equation}
The optimality system is \eqref{eq:HJB} coupled to \eqref{eq:FP} with $b_t=\grad\phi_t$.

\subsection{Hopf--Cole linearization}
Define the backward wave function
\begin{equation}
\label{eq:etastar}
\eta^\ast_t(\xv)=\exp\!\Bigl(\frac{\phi_t(\xv)}{2\gamma}\Bigr),
\qquad \grad\phi_t=2\gamma\,\frac{\grad\eta^\ast_t}{\eta^\ast_t}.
\end{equation}
A direct computation, using $\Delta\eta^\ast=\tfrac{\eta^\ast}{4\gamma^2}\|\grad\phi\|^2+\tfrac{\eta^\ast}{2\gamma}\Delta\phi$ and \eqref{eq:HJB}, gives the backward heat equation
\begin{equation}
\label{eq:heatstar}
\partial_t\eta^\ast_t=-\gamma\,\Delta\eta^\ast_t .
\end{equation}
Introduce the forward wave function through the factorization $\rho_t=\eta_t\eta^\ast_t$. With $b_t=2\gamma\grad\log\eta^\ast_t$ the Fokker--Planck equation \eqref{eq:FP} becomes
\begin{equation}
\partial_t(\eta\eta^\ast)+2\gamma\,\dive(\eta\grad\eta^\ast)=\gamma\Delta(\eta\eta^\ast);
\end{equation}
expanding and using \eqref{eq:heatstar} to cancel the terms in $\eta^\ast$ leaves the forward heat equation
\begin{equation}
\label{eq:heat}
\partial_t\eta_t=\gamma\,\Delta\eta_t .
\end{equation}
The two linear equations are coupled only through the temporal boundary
conditions (the Schr\"odinger system)
\begin{equation}
\label{eq:schrsys}
\eta_0(\xv)\,\eta^\ast_0(\xv)=\mu(\xv),\qquad
\eta_1(\xv)\,\eta^\ast_1(\xv)=\nu(\xv).
\end{equation}

\subsection{Log-potentials and IPFP}
Set $f_t=\log\eta_t$, $g_t=\log\eta^\ast_t$, so that
\begin{equation}
\rho_t=e^{f_t+g_t},\qquad b_t=2\gamma\,\grad g_t .
\end{equation}
Let $\mathcal{P}_\tau=e^{\tau\gamma\Delta}$ be the heat semigroup with homogeneous
Neumann boundary conditions. Forward and backward propagation give
$\eta_1=\mathcal{P}_1\eta_0$ and $\eta^\ast_0=\mathcal{P}_1\eta^\ast_1$; writing
$f:=\log\eta_0$ and $g:=\log\eta^\ast_1$, the system \eqref{eq:schrsys} reads
$f+\log\mathcal{P}_1[e^{g}]=\log\mu$ and $g+\log\mathcal{P}_1[e^{f}]=\log\nu$.
The iterative proportional fitting procedure (IPFP; Sinkhorn in log form) solves
these by alternation:
\begin{equation}
\label{eq:ipfp}
g^{(k+1)}=\log\nu-\log\mathcal{P}_1\!\bigl[e^{f^{(k)}}\bigr],
\qquad
f^{(k+1)}=\log\mu-\log\mathcal{P}_1\!\bigl[e^{g^{(k+1)}}\bigr].
\end{equation}
At convergence the flow is recovered by heat propagation of the endpoint
potentials, $\rho_t=\eta_t\eta^\ast_t$ and $b_t=2\gamma\grad\log\eta^\ast_t$.

% =====================================================================
\section{Reference process with a drift}
\label{sec:drift}

\subsection{Modified problem and optimality system}
Let the underlying diffusion carry a prescribed macroscopic drift
$\bt_t(\xv)$,
\begin{equation}
\label{eq:refSDE}
dX_t=\bt_t(X_t)\,dt+\sqrt{2\gamma}\,dW_t,
\end{equation}
and penalize deviation of the control from this reference (Girsanov):
\begin{equation}
\label{eq:Jdrift}
\mathcal{J}(\rho,b)=\int_0^1\!\!\int_{\Rd}\tfrac12\,\|b_t(\xv)-\bt_t(\xv)\|^2\,\rho_t(\xv)\,d\xv\,dt,
\end{equation}
under the same constraint \eqref{eq:FP}. Varying the corresponding Lagrangian in
$b$ gives the optimal control
\begin{equation}
\label{eq:boptdrift}
b_t=\bt_t+\grad\phi_t,
\end{equation}
and varying in $\rho$, followed by substitution of \eqref{eq:boptdrift}, gives the
advected Hamilton--Jacobi--Bellman equation
\begin{equation}
\label{eq:HJBdrift}
\partial_t\phi_t+\tfrac12\|\grad\phi_t\|^2+\bt_t\cdot\grad\phi_t+\gamma\,\Delta\phi_t=0 .
\end{equation}

\subsection{Transport--diffusion equations}
With the same Hopf--Cole change of variables \eqref{eq:etastar}, equation
\eqref{eq:HJBdrift} linearizes to the backward transport--diffusion equation
\begin{equation}
\label{eq:backdrift}
\partial_t\eta^\ast_t+\bt_t\cdot\grad\eta^\ast_t+\gamma\,\Delta\eta^\ast_t=0 .
\end{equation}
The forward wave function is characterized as follows.

\begin{proposition}
\label{prop:forward}
Let $\rho_t=\eta_t\eta^\ast_t$ solve \eqref{eq:FP} with optimal drift
$b_t=\bt_t+2\gamma\,\grad\eta^\ast_t/\eta^\ast_t$, and let $\eta^\ast$ satisfy
\eqref{eq:backdrift}. Then
\begin{equation}
\label{eq:forwarddrift}
\partial_t\eta_t+\dive(\bt_t\,\eta_t)-\gamma\,\Delta\eta_t=0,
\qquad\text{i.e.}\qquad
\partial_t\eta_t+\bt_t\cdot\grad\eta_t+(\dive\bt_t)\,\eta_t-\gamma\,\Delta\eta_t=0 .
\end{equation}
\end{proposition}

\begin{proof}
Insert $\rho=\eta\eta^\ast$ and $\rho b=\eta\eta^\ast\bt+2\gamma\,\eta\grad\eta^\ast$
into \eqref{eq:FP} and expand:
\begin{align*}
\partial_t(\eta\eta^\ast)&=\eta^\ast\partial_t\eta+\eta\,\partial_t\eta^\ast,\\
\dive(\eta\eta^\ast\bt)&=\eta^\ast\,\bt\cdot\grad\eta+\eta\,\bt\cdot\grad\eta^\ast+\eta\eta^\ast\,\dive\bt,\\
2\gamma\,\dive(\eta\grad\eta^\ast)&=2\gamma\bigl(\grad\eta\cdot\grad\eta^\ast+\eta\,\Delta\eta^\ast\bigr),\\
\gamma\,\Delta(\eta\eta^\ast)&=\gamma\bigl(\eta\,\Delta\eta^\ast+2\grad\eta\cdot\grad\eta^\ast+\eta^\ast\Delta\eta\bigr).
\end{align*}
Form $\partial_t\rho+\dive(\rho b)-\gamma\Delta\rho=0$ and eliminate
$\partial_t\eta^\ast$ with \eqref{eq:backdrift}, i.e.
$\eta\,\partial_t\eta^\ast=-\eta\,\bt\cdot\grad\eta^\ast-\gamma\,\eta\Delta\eta^\ast$.
The terms $\eta\,\bt\cdot\grad\eta^\ast$ cancel; the terms
$\grad\eta\cdot\grad\eta^\ast$ cancel ($+2\gamma-2\gamma$); and the terms
$\eta\Delta\eta^\ast$ cancel ($-\gamma+2\gamma-\gamma$). What remains is
$\eta^\ast\partial_t\eta+\eta^\ast\bt\cdot\grad\eta+\eta\eta^\ast\dive\bt-\gamma\eta^\ast\Delta\eta=0$;
dividing by $\eta^\ast>0$ gives \eqref{eq:forwarddrift}.
\end{proof}

\begin{remark}
The backward equation \eqref{eq:backdrift} carries the advection in
non-conservative form $\bt\cdot\grad\eta^\ast$ (the generator
$\mathcal{L}=\bt\cdot\grad+\gamma\Delta$), whereas the forward equation
\eqref{eq:forwarddrift} carries it in conservative form $\dive(\bt\,\eta)$
(the adjoint $\mathcal{L}^\ast$). The reaction term $(\dive\bt)\eta$ is therefore
present whenever $\dive\bt\neq0$ and vanishes only for a divergence-free
reference field.
\end{remark}

\subsection{Log-potentials and drifted IPFP}
With $f_t=\log\eta_t$ and $g_t=\log\eta^\ast_t$, equations
\eqref{eq:backdrift}--\eqref{eq:forwarddrift} become the coupled
Hamilton--Jacobi equations
\begin{align}
\partial_t g_t+\bt_t\cdot\grad g_t+\gamma\,\Delta g_t+\gamma\,\|\grad g_t\|^2&=0,\\
\partial_t f_t+\bt_t\cdot\grad f_t+(\dive\bt_t)-\gamma\,\Delta f_t-\gamma\,\|\grad f_t\|^2&=0.
\end{align}
Let $\mathcal{Q}_\tau=e^{\tau\mathcal{L}}$ be the transition operator of
\eqref{eq:refSDE} (the non-conservative backward semigroup) and
$\mathcal{Q}^{\dagger}_\tau=e^{\tau\mathcal{L}^\ast}$ its adjoint (the
conservative forward semigroup). Propagation of the endpoint potentials gives
$\eta^\ast_0=\mathcal{Q}_1\eta^\ast_1$ and $\eta_1=\mathcal{Q}^{\dagger}_1\eta_0$,
and the IPFP reads
\begin{equation}
\label{eq:ipfpdrift}
f^{(k+1)}=\log\mu-\log\mathcal{Q}_1\!\bigl[e^{g^{(k)}}\bigr],
\qquad
g^{(k+1)}=\log\nu-\log\mathcal{Q}^{\dagger}_1\!\bigl[e^{f^{(k+1)}}\bigr].
\end{equation}
When $\bt\equiv0$ both operators reduce to the heat semigroup $\mathcal{P}_1$ and
\eqref{eq:ipfpdrift} returns \eqref{eq:ipfp}.

\subsection{Recovery of the transport map}
The bridge yields the forward map $T$ (the conditional mean under the coupling),
\begin{equation}
\label{eq:baryT}
T(\xv)=\frac{\mathcal{Q}_1[\,\mathrm{id}\cdot e^{g}\,](\xv)}{\mathcal{Q}_1[e^{g}](\xv)}-\boldsymbol{\beta}_{\mathrm{bias}}(\xv),
\qquad
\boldsymbol{\beta}_{\mathrm{bias}}=\frac{\mathcal{Q}_1[\mathrm{id}]}{\mathcal{Q}_1[1]}-\mathrm{id},
\end{equation}
with an analogous backward map from $e^{f}$; $\boldsymbol{\beta}_{\mathrm{bias}}$
removes the identity leak of the Neumann semigroup. Evaluated in the log domain,
with a shift $c$ such that $\mathrm{id}-c>0$,
\begin{equation}
T_j(\xv)=c_j+\exp\!\Bigl(\log\mathcal{Q}_1\!\bigl[e^{\log(\mathrm{id}_j-c_j)+g}\bigr](\xv)-\log\mathcal{Q}_1\!\bigl[e^{g}\bigr](\xv)\Bigr)-(\boldsymbol{\beta}_{\mathrm{bias}})_j ,
\end{equation}
which is well defined uniformly in $\gamma$.

\subsection{Baseline transport maps (current implementation, $\bt\equiv0$)}
\label{sec:baseline}
The pipeline in its present form uses the driftless bridge, so
$\mathcal{Q}_1=\mathcal{Q}_1^{\dagger}=\mathcal{P}_1$ is the pure heat semigroup and
\eqref{eq:baryT} specialises to a \emph{pair} of barycentric maps: the forward map
$T$ (conditional mean of the target position given the source, built from $g$) and
the backward map $S$ (conditional mean of the source given the target, built from
$f$),
\begin{equation}
\label{eq:TSheat}
T(\xv)=\frac{\mathcal{P}_1[\,\mathrm{id}\,e^{g}\,](\xv)}{\mathcal{P}_1[e^{g}](\xv)}-\boldsymbol{\beta}_{\mathrm{bias}}(\xv),
\qquad
S(\xv)=\frac{\mathcal{P}_1[\,\mathrm{id}\,e^{f}\,](\xv)}{\mathcal{P}_1[e^{f}](\xv)}-\boldsymbol{\beta}_{\mathrm{bias}}(\xv),
\end{equation}
sharing the single Neumann identity-leak correction
$\boldsymbol{\beta}_{\mathrm{bias}}=\mathcal{P}_1[\mathrm{id}]/\mathcal{P}_1[1]-\mathrm{id}$.
Each of the five distinct heat solves in \eqref{eq:TSheat} is evaluated through the
stable log-domain operator $\log\mathcal{P}_1[e^{(\cdot)}]$, with the positivity
shift $\mathrm{id}\mapsto\mathrm{id}-c$ ($c=\min_\xv\mathrm{id}-1$) applied
componentwise so that $\log(\mathrm{id}-c)$ is finite; this is exactly the form
\eqref{eq:baryT} with $\mathcal{Q}_1$ replaced by $\mathcal{P}_1$. Both maps reduce
to the identity far from the transported feature (where $\boldsymbol{\beta}_{\mathrm{bias}}$
cancels the leak) and to the entropic Monge map on its support.

\subsection{Convex displacement interpolation of the Mach field}
\label{sec:cdi}
The maps $T,S$ act on the shock-feature marginals; the physical field is recovered
by a gradient-free \emph{convex displacement interpolation} (CDI). For $t\in(0,1)$,
\begin{equation}
\label{eq:cdi}
\widehat M(t,\xv)=(1-t)\,M_0\!\bigl((1-t)\,\xv+t\,S(\xv)\bigr)
                 +t\,M_1\!\bigl(t\,\xv+(1-t)\,T(\xv)\bigr),
\end{equation}
where $M_0,M_1$ are the endpoint Mach fields, sampled at the displaced pre-images by
linear interpolation on the mesh (nearest-neighbour fallback inside the solid body).
The construction is exact at the endpoints, $\widehat M(0,\cdot)=M_0$ and
$\widehat M(1,\cdot)=M_1$, and needs no map inversion: the forward pre-image
$t\,\xv+(1-t)\,T(\xv)$ pulls back $M_1$, the backward pre-image
$(1-t)\,\xv+t\,S(\xv)$ pulls back $M_0$, and the two are blended with the convex
weights $(1-t,t)$. This is the object the analytic drift of
Section~\ref{sec:diamond} is designed to improve: by pre-transporting the shocks,
$\bt$ leaves $T,S$ close to the identity and keeps \eqref{eq:cdi} well-conditioned
at small $\gamma$.

% =====================================================================
\section{Analytic reference drift for a diamond airfoil}
\label{sec:diamond}

The endpoints $\mu,\nu$ are the shock-feature marginals of the Mach fields at
freestream Mach numbers $M_0$ and $M_1$. We construct $\bt_t$ so that it
transports the shock structure between the two conditions; the bridge of
Section~\ref{sec:drift} supplies the remaining correction.

\subsection{Geometry}
The diamond is a symmetric rhombus of chord $c$, half-chord $h_c=c/2$, total
height $H=c\tan\alpha$, half-height $h_t=H/2$, centred at $(c_x,c_y)$, with
vertices
\[
\underbrace{(c_x-h_c,c_y)}_{\text{LE}},\
\underbrace{(c_x,c_y+h_t)}_{\text{upper shoulder}},\
\underbrace{(c_x+h_c,c_y)}_{\text{TE}},\
\underbrace{(c_x,c_y-h_t)}_{\text{lower shoulder}} .
\]
From the vertices, independent of $M$,
\begin{equation}
\av_{\mathrm{LE}}=(c_x-h_c,c_y),\quad
\av_{\mathrm{TE}}=(c_x+h_c,c_y),\quad
\theta_{\mathrm{LE}}=\theta_{\mathrm{TE}}=\arctan\!\frac{h_t}{h_c}=\alpha,
\end{equation}
and the shoulder produces an expansion through the angle $\Delta\theta_{\mathrm{sh}}=2\alpha$.

\subsection{Oblique-shock wave angles}
For an attached oblique shock produced by deflection $\theta$ at Mach $M$ in a gas
of ratio of specific heats $\gamma_g$, the wave angle $\beta_s$ solves the $\theta$--$\beta$--$M$ relation
\begin{equation}
\label{eq:tbm}
\tan\theta=2\,\cot\beta_s\,\frac{M^2\sin^2\beta_s-1}{M^2(\gamma_g+\cos2\beta_s)+2};
\end{equation}
denote by $\beta_s(\theta,M)$ its weak-shock root (the smaller root, exceeding
the Mach angle $\arcsin(1/M)$).

\begin{assumption}
\label{ass:attached}
$\alpha<\theta_{\max}(M):=\max_{\beta_s}\theta(\beta_s,M)$ for all
$M\in[M_0,M_1]$, so both shock families are attached, and each shock is modelled
as a straight ray from its edge.
\end{assumption}

\paragraph{Leading-edge family.} The LE shock deflects the freestream by
$\theta_{\mathrm{LE}}=\alpha$ at the freestream Mach $M$:
\begin{equation}
\beta_s^{\mathrm{LE}}(M)=\beta_s(\alpha,M).
\end{equation}

\paragraph{Trailing-edge family.} The TE shock processes the flow that has
crossed the LE shock and expanded around the shoulder. With
$\beta=\beta_s^{\mathrm{LE}}(M)$ and $M_{n1}=M\sin\beta$, the Mach behind the LE
shock is
\begin{equation}
M_{n2}^2=\frac{1+\tfrac{\gamma_g-1}{2}M_{n1}^2}{\gamma_g M_{n1}^2-\tfrac{\gamma_g-1}{2}},
\qquad
M_a=\frac{M_{n2}}{\sin(\beta-\alpha)} .
\end{equation}
The Prandtl--Meyer function
\begin{equation}
\nu(M)=\sqrt{\tfrac{\gamma_g+1}{\gamma_g-1}}\,\arctan\!\sqrt{\tfrac{\gamma_g-1}{\gamma_g+1}(M^2-1)}-\arctan\!\sqrt{M^2-1}
\end{equation}
gives the post-shoulder Mach $M_b$ through $\nu(M_b)=\nu(M_a)+2\alpha$, and the TE
wave angle follows with deflection $\theta_{\mathrm{TE}}=\alpha$ at Mach $M_b$:
\begin{equation}
\beta_s^{\mathrm{TE}}(M)=\beta_s(\alpha,M_b(M)).
\end{equation}
Both wave angles are evaluated at every interpolation node $M_i=M(t_i)$ rather
than interpolated between the endpoints.

\subsection{The drift field}
Fix the schedule $M(t)=(1-t)M_0+tM_1$ and set
$\beta_s^{\kappa}(t):=\beta_s^{\kappa}(M(t))$ for $\kappa\in\{\mathrm{LE},\mathrm{TE}\}$.
Define the ray-angle magnitudes
\begin{equation}
\psi_{\mathrm{LE}}(t)=\beta_s^{\mathrm{LE}}(t),\qquad
\psi_{\mathrm{TE}}(t)=\beta_s^{\mathrm{TE}}(t)-\alpha,
\end{equation}
the subtraction accounting for the local upstream flow along the rear face; the
lower ray angle is the negative of the upper, so the two surfaces are mirror
images. The unit ray directions and angular rates are, for $\sigma\in\{+1,-1\}$,
\begin{equation}
\dhat_{\kappa,\sigma}(t)=\bigl(\cos(\sigma\psi_\kappa),\ \sin(\sigma\psi_\kappa)\bigr),
\qquad
\omega_\kappa(t)=\dot\psi_\kappa(t)=\frac{d\beta_s^{\kappa}}{dM}\Big|_{M(t)}(M_1-M_0).
\end{equation}
Since $\beta_s$ decreases with $M$, $\omega_\kappa<0$: each shock rotates toward
the body. A point at $\xv$, riding on ray $(\kappa,\sigma)$ as it rotates about
$\av_\kappa$, has velocity $\sigma\,\omega_\kappa\,\R\,(\xv-\av_\kappa)$, using
$\R(\cos\varphi,\sin\varphi)=(-\sin\varphi,\cos\varphi)$. Localizing to the shock
by a Gaussian band of width $\sigma_w$ in the perpendicular distance, gated
downstream, define
\begin{equation}
s_{\kappa,\sigma}=(\xv-\av_\kappa)\cdot\dhat_{\kappa,\sigma},\quad
\mathrm{dist}_{\kappa,\sigma}=\bigl\|(\xv-\av_\kappa)-s_{\kappa,\sigma}\dhat_{\kappa,\sigma}\bigr\|,\quad
w_{\kappa,\sigma}=e^{-\mathrm{dist}_{\kappa,\sigma}^2/\sigma_w^2}\,\mathbf{1}[s_{\kappa,\sigma}>0].
\end{equation}
The reference drift is
\begin{equation}
\boxed{\;
\bt(t,\xv)=\sum_{\kappa\in\{\mathrm{LE},\mathrm{TE}\}}\ \sum_{\sigma\in\{+1,-1\}}
\sigma\,\omega_\kappa(t)\,w_{\kappa,\sigma}(t,\xv)\,\R\,(\xv-\av_\kappa).
\;}
\end{equation}
It vanishes in the freestream and, near each shock ray, rotates that shock about
its edge from the $M_0$ angle toward the $M_1$ angle. Because the localization
renders $\dive\bt\neq0$, the field enters the bridge through the conservative
operator $\mathcal{Q}^{\dagger}$ of \eqref{eq:ipfpdrift}, retaining the reaction
term of Proposition~\ref{prop:forward}.

\end{document}